% Part: turing-machines
% Chapter: machines-computations
% Section: configuration

\documentclass[../../../include/open-logic-section]{subfiles}

\begin{document}

\olfileid[fr]{tur}{mac}{con}
\olsection{Configurations et calculs}

\begin{explain}
Nous avons déjà suivi les configurations successives
de la machine paire. Le mécanisme imaginaire composé
d'un ruban, d'une tête de lecture-écriture et d'un
programme n'est qu'une image intuitive du calcul
d'une machine de Turing. Formellement, on peut
définir le calcul sur une entrée donnée comme une
suite de \emph{configurations}. Une configuration
comprend une suite de symboles, qui représente le
contenu du ruban à un instant du calcul, un nombre
qui indique la position de la tête et un état.
Ces données permettent de définir ce que calcule
la machine de Turing~$M$ sur une entrée donnée.
\end{explain}

\begin{defn}[Configuration]
Une \emph{configuration} de la machine de Turing
$M=\tuple{Q,\Sigma,q_0,\delta}$ est un triplet
$\tuple{C,m,q}$ tel que :
\begin{enumerate}
\item $C\in\Sigma^*$ est une suite finie de
  symboles de~$\Sigma$;
\item $m\in\Nat$ et $m<\len{C}$;
\item $q\in Q$.
\end{enumerate}
Intuitivement, la suite~$C$ représente le contenu
du ruban, depuis la case la plus à gauche jusqu'à
la dernière case non blanche ou déjà visitée;
$m$ est le numéro de la case lue par la tête
(la case la plus à gauche porte le numéro~$0$),
et $q$ est l'état courant de la machine.
\end{defn}

\begin{explain}
Une entrée possible pour la machine est une suite
de symboles, qui code généralement un nombre.
Dans la configuration initiale, le ruban porte
le marqueur de début suivi immédiatement de
l'entrée, écrite dans les cases à droite;
la tête lit la première case à droite du marqueur
et le mécanisme se trouve dans l'état initial
désigné~$q_0$.
\end{explain}

\begin{defn}[Configuration initiale]
La \emph{configuration initiale} de~$M$ pour
l'entrée $I\in\Sigma^*$ est
\[
\tuple{\TMendtape \frown I \frown \TMblank, 1, q_0}.
\]
La case blanche finale assure que la position~$1$
existe aussi lorsque l'entrée est vide.\footnote{Correction éditoriale : la suite source n'ajoute aucune case blanche. Pour une entrée vide, elle aurait une seule case et la position $1$ violerait $m<\len{C}$. Nous représentons donc explicitement une case blanche supplémentaire.}
\end{defn}

\begin{explain}
Le symbole~$\frown$ désigne la \emph{concaténation} :
la chaîne d'entrée commence immédiatement à droite
du marqueur de début du ruban. La case blanche
supplémentaire ne fait que prolonger la représentation
finie du ruban, qui est blanc au-delà de cette suite.
\end{explain}

\begin{defn}
On dit qu'une configuration $\tuple{C,m,q}$
\emph{produit en une étape la configuration
$\tuple{C',m',q'}$} (selon~$M$) si et seulement si :
\begin{enumerate}
\item le symbole de rang~$m$ dans~$C$ est~$\sigma$;
\item les instructions de~$M$ prescrivent
  $\delta(q,\sigma)=\tuple{q',\sigma',D}$;
\item le symbole de rang~$m$ dans~$C'$ est~$\sigma'$;
\item
\begin{enumerate}
\item soit $D=L$ et $m'=m-1$ si $m>0$,
  et $m'=0$ sinon;
\item soit $D=R$ et $m'=m+1$;
\item soit $D=N$ et $m'=m$;
\end{enumerate}
\item si $m'=\len{C}$, alors
  $\len{C'}=\len{C}+1$ et le symbole de rang~$m'$
  dans~$C'$ est~$\TMblank$; sinon,
  $\len{C'}=\len{C}$;
\item pour tout~$i$ tel que $i<\len{C}$ et
  $i\neq m$, on a $C'(i)=C(i)$.
\end{enumerate}
\end{defn}

\begin{defn}\ollabel{defn:run-output}
Une \emph{exécution de~$M$ sur l'entrée~$I$}
est une suite, finie ou infinie, de configurations
$C_i$ de~$M$, dont $C_0$ est la configuration
initiale pour~$I$ et dont chaque paire de
configurations consécutives $C_i,C_{i+1}$ est
liée par une étape.

On dit que~$M$ \emph{s'arrête sur l'entrée~$I$
après $k$ étapes} si $C_k=\tuple{C,m,q}$,
si le symbole de rang~$m$ dans~$C$ est~$\sigma$
et si $\delta(q,\sigma)$ n'est pas définie.
Dans ce cas, la \emph{sortie} de~$M$ sur~$I$
est la chaîne~$O$ ne se terminant pas par
$\TMblank$ et telle que
$C=C(0)\concat O\concat\TMblank^j$ pour un
certain~$j\in\Nat$. Autrement dit, la première
case est écartée, qu'elle contienne encore
$\TMendtape$ ou que la machine l'ait écrasé.\footnote{Correction éditoriale : la source exige $C=\TMendtape\concat O\concat\TMblank^j$, alors qu'elle autorise l'écrasement du marqueur. Nous utilisons le contenu actuel $C(0)$ pour définir la sortie de toute machine arrêtée.}
La notation $\TMblank^j$ désigne une suite de
$j$ symboles $\TMblank$.
\end{defn}

\begin{explain}
Selon cette définition, la sortie~$O$ de~$M$
se termine toujours par un symbole autre que
$\TMblank$; si, à l'étape~$k$, toutes les cases
représentées après la première sont blanches,
$O$ est la chaîne vide.
\end{explain}

\end{document}
